\section{Fazit und Ausblick}

Ziel dieser Arbeit war die Untersuchung eines semi-supervisierten Ansatzes zur
automatischen Annotation von UAV-Videodaten für die Detektion potenzieller
Brutstätten von \textit{Aedes aegypti}. Im Mittelpunkt stand die Frage, ob eine
begrenzte Menge manuell annotierter Frames ausreicht, um ein YOLO-basiertes
Objekterkennungsmodell zu trainieren und dieses anschließend zur automatischen
Erweiterung des Trainingsdatensatzes einzusetzen. Zusätzlich wurde untersucht,
inwieweit die zeitliche Struktur von Videodaten durch tracking-basierte
Label-Propagation genutzt werden kann, um automatisch erzeugte Annotationen
zu stabilisieren und als zusätzliche Trainingsdaten zu verwenden.

Als Ausgangspunkt wurde ein Baseline-Modell auf manuell annotierten Frames des
MBG-V2-Datensatzes trainiert. Dieses Modell diente als Referenz für die
nachfolgenden semi-supervisierten Experimente. Anschließend wurden mithilfe des
Baseline-Modells zusätzliche Pseudo-Labels auf nicht manuell annotierten Frames
erzeugt und in einem erweiterten Trainingsdatensatz verwendet. In einem weiteren
Schritt wurde ByteTrack eingesetzt, um Detektionen über mehrere
aufeinanderfolgende Frames hinweg zu verknüpfen. Die daraus gewonnenen
tracking-basierten Pseudo-Labels wurden gefiltert, mit den manuellen
Trainingsdaten kombiniert und in einem eigenen Retraining-Schritt quantitativ
evaluiert.

Die Ergebnisse zeigen, dass das ausschließlich auf manuellen Annotationen
trainierte Baseline-Modell im finalen Vergleich die beste Gesamtleistung erreicht.
Auf demselben manuell annotierten Validierungssplit erzielt es eine Precision von
0{,}831, einen Recall von 0{,}485, einen mAP@50-Wert von 0{,}572 und einen
mAP@50--95-Wert von 0{,}343. Besonders häufige und visuell klar erkennbare
Klassen wie \textit{pool} und \textit{watertank} konnten vergleichsweise zuverlässig
erkannt werden. Gleichzeitig zeigten sich deutliche Einschränkungen bei kleinen,
seltenen oder visuell schwer unterscheidbaren Klassen. Diese Objekte wurden
häufig mit dem Hintergrund verwechselt oder gar nicht erkannt.

Die detection-basierte Pseudo-Labeling-Stufe führte zu einer deutlichen
Erweiterung der Datenbasis. Insbesondere mehrere zuvor unterrepräsentierte
Klassen wie \textit{bucket}, \textit{tire}, \textit{large trash bin} und
\textit{plastic bag} erhielten zusätzliche Trainingsinstanzen. Die reine Erhöhung
der Label-Anzahl führte jedoch nicht automatisch zu einer Verbesserung der
Modellleistung. Im finalen Vergleich erreichte das detection-basierte
Pseudo-Labeling eine Precision von 0{,}692, einen Recall von 0{,}489, einen
mAP@50-Wert von 0{,}564 und einen mAP@50--95-Wert von 0{,}323. Der Recall
lag damit minimal über dem Baseline-Modell, während die Precision deutlich sank.
Dieses Ergebnis zeigt, dass automatisch erzeugte Labels zwar die Klassenabdeckung
verbessern können, gleichzeitig aber auch Rauschen und Fehlannotationen in den
Trainingsprozess einbringen.

Die tracking-basierte Label-Propagation erzeugte konservativere Pseudo-Labels.
Das entsprechend trainierte Modell erreichte eine Precision von 0{,}765, einen
Recall von 0{,}445, einen mAP@50-Wert von 0{,}543 und einen mAP@50--95-Wert
von 0{,}327. Damit übertraf es das detection-basierte Pseudo-Labeling hinsichtlich
Precision, blieb jedoch unterhalb des Baseline-Modells. Die Ergebnisse zeigen,
dass Tracking zur Qualitätskontrolle automatisch erzeugter Labels beitragen kann,
indem kurzlebige oder inkonsistente Einzelvorhersagen stärker herausgefiltert
werden. Gleichzeitig reicht die tracking-basierte Filterung in der vorliegenden
Konfiguration nicht aus, um die Leistung des manuell trainierten Baseline-Modells
zu übertreffen.

Ein zentrales Ergebnis dieser Arbeit ist daher, dass die Qualität der Pseudo-Labels
entscheidender ist als deren reine Anzahl. Wenn Fehlvorhersagen des
Baseline-Modells ungefiltert oder nur schwach gefiltert in den erweiterten
Trainingsdatensatz übernommen werden, können sich diese Fehler im
anschließenden Training verstärken. Besonders sichtbar wurde dies bei False
Positives für die Klasse \textit{watertank} sowie bei schwierigen
Hintergrundstrukturen, die potenziellen Brutstätten visuell ähneln. Die Ergebnisse
verdeutlichen damit eine typische Herausforderung semi-supervisierter
Objekterkennung: Automatische Annotation kann den manuellen Aufwand
reduzieren, erfordert jedoch geeignete Mechanismen zur Qualitätssicherung.

Zusammenfassend zeigt diese Arbeit, dass semi-supervisierte Annotation mit
Pseudo-Labeling und tracking-basierter Label-Propagation grundsätzlich geeignet
ist, zusätzliche Trainingsdaten aus UAV-Videosequenzen zu erschließen. Die
untersuchten Verfahren konnten die Datenbasis erweitern und typische Probleme
automatischer Annotation sichtbar machen. Eine klare Leistungssteigerung
gegenüber dem manuell trainierten Baseline-Modell wurde jedoch nicht erreicht.
Tracking erwies sich dabei als hilfreicher Mechanismus zur konservativeren
Auswahl von Pseudo-Labels, ersetzt aber keine sorgfältige Filterung, Validierung
und gegebenenfalls manuelle Kontrolle automatisch erzeugter Annotationen.

\subsection{Beantwortung der Forschungsfragen}

Die erste Forschungsfrage bezog sich darauf, inwieweit ein YOLO-basiertes
Baseline-Modell auf Grundlage einer begrenzten Menge manuell annotierter
UAV-Frames potenzielle Brutstätten erkennen kann. Die Ergebnisse zeigen, dass
dies grundsätzlich möglich ist. Das Baseline-Modell erreicht im finalen Vergleich
eine Precision von 0{,}831, einen Recall von 0{,}485, einen mAP@50-Wert von
0{,}572 und einen mAP@50--95-Wert von 0{,}343. Besonders häufige und visuell
eindeutige Klassen wie \textit{pool} und \textit{watertank} werden zuverlässig
erkannt. Gleichzeitig bleibt der Recall begrenzt, da kleine und seltene Objekte
häufig übersehen werden.

Die zweite Forschungsfrage untersuchte, ob automatische Pseudo-Label-Generierung
den Trainingsdatensatz sinnvoll erweitern und die Klassenabdeckung verbessern
kann. Diese Frage kann teilweise positiv beantwortet werden. Die Pseudo-Labels
erweiterten die Datenbasis deutlich und reduzierten die extreme Ungleichverteilung
mehrerer Klassen. Besonders unterrepräsentierte Klassen erhielten zusätzliche
Beispiele. Allerdings war die Qualität dieser Labels nicht in allen Fällen
ausreichend, um eine Verbesserung der Modellleistung zu erreichen.

Die dritte Forschungsfrage betraf die Auswirkungen der Pseudo-Labels auf
Precision, Recall, mAP@50 und mAP@50--95. Die Ergebnisse zeigen, dass die
detection-basierte Pseudo-Labeling-Stufe trotz größerer Datenmenge keine klare
Leistungssteigerung erzielte. Zwar stieg der Recall im finalen Vergleich leicht von
0{,}485 auf 0{,}489, gleichzeitig sank die Precision jedoch deutlich von 0{,}831
auf 0{,}692. Auch der mAP@50--95-Wert lag mit 0{,}323 unter dem Wert des
Baseline-Modells. Dies deutet darauf hin, dass die zusätzlichen Labels teilweise
fehlerhafte oder unsichere Annotationen enthielten und dadurch den Trainingsprozess
nicht eindeutig verbesserten.

Die vierte Forschungsfrage bezog sich auf die tracking-basierte Label-Propagation.
Die Ergebnisse zeigen, dass Tracking einen sinnvollen Beitrag zur zeitlichen
Stabilisierung und Qualitätskontrolle automatisch erzeugter Detektionen leisten
kann. Durch die Nutzung stabiler Objekttracks konnten konservativere
Pseudo-Labels erzeugt werden. Das tracking-basierte Modell erreichte eine
Precision von 0{,}765 und lag damit deutlich über dem detection-basierten
Pseudo-Labeling. Gleichzeitig blieb es mit einem Recall von 0{,}445, einem
mAP@50-Wert von 0{,}543 und einem mAP@50--95-Wert von 0{,}327 unterhalb
des Baseline-Modells. Tracking verbessert somit die Zuverlässigkeit der
Pseudo-Label-Auswahl, führte in dieser Arbeit jedoch nicht zu einer höheren
Gesamtleistung als das manuell trainierte Modell.

Die fünfte Forschungsfrage untersuchte typische Fehlerquellen bei der
automatischen Annotation von UAV-Videodaten. Als zentrale Fehlerquellen wurden
kleine Objektgrößen, komplexe Hintergrundstrukturen, Klassenungleichgewicht,
visuell ähnliche Objekte und instabile Detektionen identifiziert. Besonders
problematisch sind Objekte, die nur wenige Pixel groß sind oder sich farblich und
strukturell kaum vom Hintergrund unterscheiden. Die Fehleranalyse zeigt, dass
viele Probleme nicht primär durch falsche Klassenzuordnungen zwischen
Objektklassen entstehen, sondern durch nicht erkannte Objekte und Verwechslungen
mit dem Hintergrund.

\subsection{Limitationen}

Die Ergebnisse dieser Arbeit müssen im Kontext mehrerer Einschränkungen
betrachtet werden. Erstens basiert die Auswertung auf einer begrenzten Menge
manuell annotierter Daten. Obwohl der MBG-V2-Datensatz umfangreiche
UAV-Videodaten enthält, sind nicht alle Frames vollständig manuell annotiert.
Dadurch bleibt die verfügbare Menge zuverlässiger Ground-Truth-Daten begrenzt.

Zweitens ist die Klassenverteilung stark unausgewogen. Die Klasse
\textit{watertank} ist deutlich häufiger vertreten als mehrere andere Klassen. Diese
Ungleichverteilung beeinflusst sowohl das Training als auch die Evaluation.
Häufige Klassen werden stabiler gelernt, während seltene Klassen nur eingeschränkt
bewertet werden können. Besonders die Ergebnisse für Klassen mit sehr wenigen
Validierungsinstanzen, etwa \textit{bottle}, \textit{plastic bag} oder
\textit{large trash bin}, müssen daher vorsichtig interpretiert werden.

Drittens wurden die Pseudo-Labels in dieser Arbeit mit vergleichsweise einfachen
Filterstrategien erzeugt. Eine feinere Qualitätskontrolle, beispielsweise durch
klassenabhängige Confidence-Schwellenwerte, automatische Unsicherheitsmaße oder
eine systematische manuelle Stichprobenprüfung, wurde nicht vollständig
ausgeschöpft. Dadurch konnten fehlerhafte Pseudo-Labels in den erweiterten
Trainingsdatensatz gelangen.

Viertens konnte die tracking-basierte Label-Propagation das Baseline-Modell nicht
übertreffen. Obwohl sie quantitativ durch einen eigenen Retraining-Schritt
ausgewertet wurde und eine höhere Precision als das detection-basierte
Pseudo-Labeling erreichte, blieb die Gesamtleistung unterhalb des manuell
trainierten Modells. Dies zeigt, dass zeitliche Konsistenz allein nicht ausreicht,
wenn die zugrunde liegenden Detektionen unsicher sind oder kleine Objekte bereits
vom Baseline-Modell häufig übersehen werden.

Fünftens beschränkt sich die Arbeit auf eine ausgewählte Trainings- und
Filterkonfiguration. Andere YOLO-Modellgrößen, höhere Eingangsauflösungen,
Patch-basierte Verfahren oder alternative Tracking-Algorithmen wurden nicht
systematisch verglichen. Die Ergebnisse sind daher als experimentelle Bewertung
des vorgeschlagenen Workflows zu verstehen und nicht als abschließende
Optimierung aller möglichen Modell- und Filtervarianten.

\subsection{Ausblick}

Aufbauend auf den Ergebnissen dieser Arbeit ergeben sich mehrere Ansätze für
zukünftige Arbeiten. Ein wichtiger nächster Schritt besteht in einer stärkeren
Qualitätskontrolle der Pseudo-Labels. Dazu könnten höhere oder klassenabhängige
Confidence-Schwellenwerte eingesetzt werden. Für häufige Klassen wie
\textit{watertank} könnte ein strengerer Schwellenwert sinnvoll sein, um False
Positives zu reduzieren. Für seltene Klassen könnten dagegen angepasste Strategien
erforderlich sein, um genügend zusätzliche Trainingsbeispiele zu erhalten, ohne zu
viele Fehlannotationen zu übernehmen.

Ein weiterer Ansatz ist die systematische Weiterentwicklung tracking-basierter
Filterkriterien. In dieser Arbeit wurden Tracklänge, mittlere Konfidenz und
Klassenkonsistenz zur Auswahl tracking-basierter Pseudo-Labels genutzt. Zukünftige
Arbeiten könnten diese Kriterien erweitern, beispielsweise durch räumliche
Stabilitätsmaße, Bewegungsmodelle, Track-Unterbrechungen oder Unsicherheitsmaße
über mehrere Frames hinweg. Dadurch ließen sich kurzlebige Einzelvorhersagen
und instabile False Positives möglicherweise noch zuverlässiger reduzieren.

Darüber hinaus könnte der Datensatz durch zusätzliche manuell annotierte
Videosequenzen erweitert werden. Besonders für seltene und kleine Klassen wäre
eine gezielte Erweiterung sinnvoll. Eine größere und ausgewogenere Datenbasis
könnte die Erkennung schwieriger Klassen verbessern und gleichzeitig die Qualität
der automatisch erzeugten Labels erhöhen.

Auch eine weiterführende Modelloptimierung wäre denkbar. Dazu gehören
alternative YOLO-Modellgrößen, angepasste Bildauflösungen, Patch-basierte
Verarbeitung großer UAV-Bilder oder spezialisierte Augmentierungsstrategien für
kleine Objekte. Gerade bei UAV-Aufnahmen kann eine höhere Eingangsauflösung
oder eine Aufteilung der Bilder in kleinere Ausschnitte helfen, kleine Objekte
besser sichtbar zu machen.

Zusätzlich könnten zukünftige Arbeiten die tracking-basierte Label-Propagation
mit aktiven Lernstrategien kombinieren. Dabei könnten besonders unsichere oder
widersprüchliche Tracks gezielt für eine manuelle Nachkontrolle ausgewählt
werden. Auf diese Weise ließe sich der manuelle Annotierungsaufwand weiterhin
begrenzen, während die Qualität der Trainingsdaten gezielt verbessert wird.

Zusammenfassend zeigt diese Arbeit, dass semi-supervisierte Annotation mit
Pseudo-Labeling und Tracking ein vielversprechender Ansatz zur Reduktion
manuellen Annotierungsaufwands in UAV-Videodaten ist. Gleichzeitig machen die
Ergebnisse deutlich, dass automatische Labels nicht unkritisch übernommen werden
dürfen. Entscheidend für den Erfolg des Ansatzes ist eine kontrollierte Auswahl
qualitativ hochwertiger Pseudo-Labels. Tracking-basierte zeitliche Konsistenz
stellt hierfür einen sinnvollen ergänzenden Mechanismus dar und bietet eine
wichtige Grundlage für zukünftige Verbesserungen automatischer
Annotationsverfahren in videobasierten UAV-Datensätzen.